\documentclass{article}
\usepackage[utf8]{inputenc}
\usepackage{amsmath}
\usepackage{geometry}
\geometry{margin=2cm}
\begin{document}

\section*{Análise da Questão 162 — ENEM 2024 — Matemática}

\subsection*{Enunciado}
Uma imobiliária oferece financiamento de apartamentos com 100 parcelas mensais, onde a primeira parcela é paga no momento da compra com valor fixo $P$. Cada parcela subsequente é reajustada por um acréscimo percentual fixo formando uma progressão geométrica com razão $r>1$. A imobiliária pagará todas as parcelas correspondentes ao mês de aniversário do comprador. Um cliente com aniversário em maio deseja escolher o mês para comprar o apartamento reduzindo o valor total pago. Qual mês ele deve escolher?

\subsection*{Solução Técnica}

Para determinar o mês ideal para a compra do apartamento, consideramos a sequência das 100 parcelas como uma progressão geométrica crescente. Seja $P$ o valor da primeira parcela e $r > 1$ a razão da progressão.

A n-ésima parcela é dada por:
\[
P_n = P \cdot r^{n-1}
\]

Como o comprador faz aniversário em maio, parcelas referentes a maio são pagas pela imobiliária.

O total pago pelo comprador é:
\[
S = \sum_{n=1}^{100} P \cdot r^{n-1} - \sum_{k \in M} P \cdot r^{k-1}
\]
onde $M$ é o conjunto de índices das parcelas em maio.

Comparando os totais para os meses candidatos (fevereiro, abril, maio, junho, agosto), verifica-se que iniciar em \textbf{fevereiro} minimiza o valor total pago.

\subsection*{Explicação Didática}

\textbf{Resumo:} O valor das parcelas aumenta mês a mês, enquanto as parcelas referentes ao mês de aniversário (maio) são gratuitas para o comprador. Analisando a sequência e as parcelas descontadas, identifica-se que comprar em fevereiro é mais vantajoso.

\textbf{Passos para compreensão:}
\begin{enumerate}
  \item Identificar que as parcelas formam uma progressão geométrica crescente.
  \item Mapear quais parcelas cairão em maio conforme o mês de início.
  \item Calcular o somatório dos pagamentos efetuados e dos descontos.
  \item Comparar os totais pagos para cada opção.
  \item Concluir que fevereiro gera o menor custo.
\end{enumerate}

\subsection*{Erros Comuns}
\begin{itemize}
  \item Supor que comprar no mês do aniversário (maio) é sempre melhor.
  \item Não considerar o crescimento progressivo dos valores das parcelas.
  \item Contar incorretamente as parcelas isentas.
\end{itemize}

\subsection*{Dicas para Ensino}
\begin{itemize}
  \item Usar exemplos numéricos simples.
  \item Fazer uso de calendário para visualização.
  \item Apresentar a progressão geométrica de forma intuitiva.
  \item Praticar com diferentes meses de aniversário.
\end{itemize}

\subsection*{Análise dos Distratores}
\begin{itemize}
  \item \textbf{Abril:} Estratégia errada por não maximizar a isenção das parcelas.
  \item \textbf{Maio:} Intuitivo, mas não otimiza o total pago.
  \item \textbf{Junho:} Falsa impressão de mais parcelas gratuitas.
  \item \textbf{Agosto:} Desconsidera o impacto da progressão geométrica e aumenta o custo.
\end{itemize}

\subsection*{Perfil da Questão}
\begin{itemize}
  \item Habilidade Primária: Raciocínio Matemático e Álgebra.
  \item Habilidades Secundárias: Progressão Geométrica, Análise de Sequências, Raciocínio Financeiro.
  \item Demanda Cognitiva: Médio.
  \item Dificuldade Estimada: Média.
  \item Observações: Requer modelagem e análise combinada de sequência e calendários.
\end{itemize}

\subsection*{Conclusão}
A análise mostra que a resposta correta é \textbf{fevereiro}, que minimiza o valor total pago pelo comprador devido ao alinhamento dos meses das parcelas gratuitas e a progressão geométrica do acréscimo das parcelas.

\end{document}